% 第6章 6.4.2：正权支撑分量与零模。顶点 1-2-3 为单位权路径，4-5 为单位权边，
% 3-4 为零权结构边（虚线），6 为孤立点。正权支撑分量 C1={1,2,3}，C2={4,5}，C3={6}。
% 加权度 d=(1,2,1,1,1,0)。L 的零空间基：1_{C1},1_{C2},1_{C3}；
% L_sym 的零空间基：D^{1/2}1_{C1}=(1,sqrt2,1,0,0,0)，D^{1/2}1_{C2}=(0,0,0,1,1,0)，e_6（孤点坐标方向）。
\begin{tikzpicture}[x=1mm,y=1mm,font=\figurefont\footnotesize,text=FigureInk,
  vertex/.style={circle,draw=FigureStroke,fill=white,line width=.8pt,minimum size=6mm,inner sep=0pt},
  edge/.style={draw=FigureStroke,line width=.8pt},
  lab/.style={anchor=east,font=\figurefont\footnotesize},
  num/.style={font=\figurefont\footnotesize},
  note/.style={font=\FigureNoteFont\fontsize{7.5}{9}\selectfont,text=FigureMuted}]
\fill[FigureBlueFill,rounded corners=3mm] (23,41) rectangle (61,51);
\fill[FigureGreenFill,rounded corners=3mm] (67,41) rectangle (91,51);
\fill[FigureYellowFill,rounded corners=3mm] (97,41) rectangle (107,51);
\foreach \i/\x in {1/28,2/42,3/56,4/72,5/86,6/102}{\node[vertex] (v\i) at (\x,46) {$\i$};}
\draw[edge] (v1)--(v2)--(v3);
\draw[edge] (v4)--(v5);
\draw[draw=FigureMuted,line width=.8pt,dashed] (v3)--(v4) node[midway,above=2.4mm,note]{$w=0$};
\node[note] at (42,54) {$C_1$};
\node[note] at (79,54) {$C_2$};
\node[note] at (102,54) {$C_3$};
\node[lab] at (20,34) {度$d_i$};
\foreach \x/\v in {28/1,42/2,56/1,72/1,86/1,102/0}{\node[num,text=FigureMuted] at (\x,34) {$\v$};}
\draw[FigureGrid,line width=.4pt] (0,30.5)--(107,30.5);
\node[anchor=west] at (0,27.5) {$\bm L$的零模};
\foreach \y/\name/\vals in {23/{\bm 1_{C_1}}/{1,1,1,0,0,0},18/{\bm 1_{C_2}}/{0,0,0,1,1,0},13/{\bm 1_{C_3}}/{0,0,0,0,0,1}}{
  \node[lab] at (20,\y) {$\name$};
  \foreach \v [count=\k] in \vals {
    \pgfmathsetmacro{\xx}{{28,42,56,72,86,102}[\k-1]}
    \node[num,text=FigureBlue] at (\xx,\y) {$\v$};}}
\draw[FigureGrid,line width=.4pt] (0,8.5)--(107,8.5);
\node[anchor=west] at (0,5.5) {$\bm L_{\mathrm{sym}}$的零模};
\foreach \y/\name/\vals in {1/{\bm D^{1/2}\bm 1_{C_1}}/{1,\sqrt2,1,0,0,0},-4.5/{\bm D^{1/2}\bm 1_{C_2}}/{0,0,0,1,1,0},-10/{\bm e_6}/{0,0,0,0,0,1}}{
  \node[lab] at (20,\y) {$\name$};
  \foreach \v [count=\k] in \vals {
    \pgfmathsetmacro{\xx}{{28,42,56,72,86,102}[\k-1]}
    \node[num,text=FigureGreen] at (\xx,\y) {$\v$};}}
\end{tikzpicture}
